\section{The current proof state: an honest audit}
\label{sec:proof-state}

This section is the corrected account of the program's standing. It
states and proves the three corrections to the UC13--UC14 ledger,
isolates $Y_1$ as a precisely named open problem, and reports the state
of the machine-checked formalisation. It is written so that a future
reader does not re-propagate the known defects.

\subsection{$Y_2$: the obstruction-class landing is mechanism-sound}
\label{ssec:y2}

Recall $Y_2$ (Definition~\ref{def:y1-y2}): the obstruction class lands
in the level-$1$ Walsh isotypes. The ledger carried, at one point,
specific evidence that $Y_2$ \emph{might be false} --- the natural
machine encoding placed the obstruction on a cube vertex and proved it
nonzero. An honest, neutral machine-checked formalisation resolves the
question. The result is that $Y_2$'s \emph{mechanism} is sound and the
``evidence against'' was an encoding artefact.

\begin{theorem}[$Y_2$ mechanism, machine-checked]
\label{thm:y2-mechanism}
The following hold and have been verified in Lean against
\texttt{mathlib} \textup(no \texttt{sorry}, no project axiom; the proofs
depend only on the three \texttt{mathlib} foundational axioms\textup).
\begin{enumerate}[label=\textup{(\arabic*)},leftmargin=2.4em]
\item \textup{(landing mechanism)} A $(\Z/2)^n$-equivariant additive
  $\Q$-linear map preserves the Walsh isotype of every isotypic
  element. \textup(This is
  Lemma~\ref{lem:equivariant-preserves}, machine-checked.\textup)
\item \textup{(no additive level shift)} Multiplication by a Walsh
  character $\walsh U$ shifts $\walsh S\mapsto\walsh{S\bigtriangleup
  U}$ and is \emph{not} $(\Z/2)^n$-equivariant; hence no additive
  equivariant map shifts Walsh level
  \textup(Lemma~\ref{lem:cup-shifts}\textup). The competing UC11~\S5.3
  claim --- that the additive \v{C}ech edge map increments Walsh level
  --- is therefore representation-theoretically \emph{impossible}.
\item \textup{(encoding diagnosis)} A single cube vertex $\delta_A$ is
  in \emph{no} Walsh isotype; by Walsh inversion it has a nonzero
  component of equal magnitude in \emph{every} one of the $2^n$
  isotypes. The machine encoding \texttt{obstructionClass} used in the
  formalisation is such a single vertex, hence isotype-blind: it
  supports neither $Y_2$ nor its negation.
\item \textup{(label correction)} The Lean lemma
  \texttt{topVertex\_not\_coboundary} factors through the augmentation,
  which is the \emph{trivial} $\walsh\emptyset$ coefficient at degree
  $0$; it detects the $H^0$ connectedness generator, not the top
  $\walsh{[n]}$ isotype at degree $n-1$. The earlier reading that it
  ``proved the opposite of $Y_2$'' rested on mislabelling this
  $\walsh\emptyset$/degree-$0$ statement as a $\walsh{[n]}$/degree-%
  $(n-1)$ one.
\end{enumerate}
\end{theorem}

\begin{status}
$Y_2$ is \GREEN{mechanism-sound}. Items (1)--(2) establish that the
landing argument of UC13~\S2.4.1 uses a correct mechanism and that the
abandoned UC11~\S5.3 alternative is impossible; items (3)--(4) show the
``formalisation contradicts $Y_2$'' alarm was an artefact of a
non-faithful single-vertex encoding plus an isotype mislabelling. One
genuine caveat remains: $Y_2$ as a \emph{complete} statement is
``mechanism $+$ the source data of $\fobs$ being genuinely level-$1$
Walsh''. The source-data half is a design property of $\fobs$ --- the
bias is, up to sign, the level-$1$ $\walsh{\{x\}}$-coefficient of the
family indicator (Proposition~\ref{prop:four-forms}(O3)) --- and the
remaining obligation is the structural check that the \v{C}ech double
complex of $\fobs$ carries no cup product
(Remark~\ref{rem:source-level1}). That is a checkable structural
property, not a coin-flip. $Y_2$ is the \emph{sound} half of the
two-part contradiction.
\end{status}

\subsection{$Y_1$: the level-$1$ Walsh vanishing is open}
\label{ssec:y1}

$Y_1$ (Definition~\ref{def:y1-y2}) is the vanishing half. We first
record why the circulated proof is invalid, then the auxiliary lemma
that was supposed to repair it and why it is false, then the valid
reduction that isolates the genuine open problem.

\subsubsection{The trace-exactness error}

\begin{correction}[UC13 \S4.5 conflates trace-exactness with vanishing]
\label{corr:trace-exact}
The proof of $Y_1$ given in UC13~\S4.5 (Theorem~4.5.1) is logically
invalid. It is structured as follows: for $S\subsetneq[n]$ and
$x\in[n]\setminus S$, a twisted-bridge chain-homotopy identity (the
$\walsh{\{x\}}$-conjugate of Construction~\ref{constr:bridge}) shows
that the trace projection of any $\walsh S$-cocycle $\omega$ is a
coboundary,
\[
  \iota_x^{\mathcal D\,*}\omega
  \;=\;\pm\tfrac12\,d\bigl(h_x^{\mathrm{tw}}\omega\bigr);
\]
\S4.5 then iterates this across all $n-|S|$ coordinates outside $S$ and
concludes ``each class is exact at every level'', hence
$\isotype Sk(\Xn)=0$.

\emph{The defect.} The identity says the \emph{induced map}
$\iota_x^{\mathcal D\,*}$ is the \emph{zero map} on
$\walsh S$-cohomology --- a statement about the \emph{image} of a map
out of $\isotype Sk(\Xn)$. It says nothing about whether the
\emph{domain} $\isotype Sk(\Xn)$ is zero: a zero map out of a nonzero
group is perfectly possible. Composing zero maps down the iteration
yields, again, only that each trace is zero, never that $\omega$
itself is a coboundary. To convert ``$\iota_x^{\mathcal D\,*}=0$ on
cohomology'' into ``$\isotype Sk(\Xn)=0$'' one needs
$\iota_x^{\mathcal D\,*}$ to be \emph{injective}; \S4.5's parenthetical
justification of injectivity (``at most isomorphic by the same argument
at $n-1$'') is circular. The gap is real and unfilled. This is the
same ``trace exact $\neq$ class zero'' distinction flagged in
Remark~\ref{rem:bridge-vs-y1}; the sibling top-Walsh argument
(UC13~\S5.2) was honest about it (``we need a more careful argument''),
\S4.5 was not.
\end{correction}

\subsubsection{The false repair}

\begin{correction}[UC14 ``$R2$'' is provably false]
\label{corr:r2}
UC14~$R2$ (Theorems~2.6.2 and~2.7.1) is the lemma advertised as
repairing Correction~\ref{corr:trace-exact}'s gap. It upgrades the
cohomological \emph{surjection} of UC13 Lemma~4.4.1 to a chain-level
\emph{isomorphism} on the $\walsh S$-isotype,
\[
  C^*(\Xn)_{\walsh S}\;=\;C^*\bigl(X(\Dx)\bigr)_{\walsh S}
  \qquad(S\subsetneq[n],\ x\notin S),
\]
which would make $\iota_x^{\mathcal D\,*}$ an isomorphism, hence
injective, closing \S4.5. \textbf{The isomorphism is false.}

\emph{The refutation.} Apply $R2$ at $S=[n]\setminus\{x\}$ (a legal
instance: $R2$ is asserted for every $S\subsetneq[n]$) and $k=n-2$.
The matched subcategory $\Dx$ has $X(\Dx)\cong\Xnm\times I$, a
cylinder; cohomology is homotopy-invariant, and $[n]\setminus\{x\}$ is
the \emph{full} ground set of $\Xnm$, so $\walsh{[n]\setminus\{x\}}$ is
its \emph{top} character. Hence
\[
  \widetilde H^{n-2}\bigl(X(\Dx)\bigr)_{\walsh{[n]\setminus\{x\}}}
  \;\cong\;V^{n-2}_{[n-1]}(X^{\cap}_{n-1})
  \;\cong\;\sgn_{\Sym_{n-1}}\;\neq\;0,
\]
the program's own sphere class --- the very class
Theorem~\ref{thm:trace-injectivity} and UC14~$R3$ use as a load-bearing
\emph{nonzero} input. But the target $\isotype{[n]\setminus\{x\}}{n-2}
(\Xn)$ is a \emph{lower} isotype, which $Y_1$ (and
Conjecture~\ref{conj:uc10-1}) require to vanish. A chain-level
isomorphism cannot carry a nonzero group onto a zero group. So $R2$
contradicts both $Y_1$ and its own sibling $R3$, inside one document;
the contradiction is concrete already at $n=3$.

\emph{The root cause.} $R2$'s single-family, cell-level identification
$C^*(X(\F))_{\walsh S}=C^*(X(\Dx\F))_{\walsh S}$ is plausibly correct.
The fatal step is the \emph{hocolim lift} (Theorem~2.6.2): it claims a
per-object identification on the subcategory $\Dx$ lifts to a chain
isomorphism of the homotopy colimits. It does not. A hocolim is a
Bousfield--Kan complex summed over bar strings through the
\emph{entire} category $\Cn$; UC14 Lemma~2.6.1 establishes only that
$\Dx$ is a categorical \emph{retraction}, which makes $X(\Dx)$ a
\emph{direct summand} of $\Xn$, not an isomorph. The general principle:
\textbf{termwise \emph{vanishing} lifts through a hocolim for free;
termwise \emph{isomorphism} does not.} Any argument shaped like
``identify family-by-family on a subcategory, conclude the hocolims
agree'' is suspect, and should be checked against this.
\end{correction}

\begin{remark}
\label{rem:r2-lesson}
Correction~\ref{corr:r2} retracts the internal ``\GREEN{$R1+R2+R3$ all
closed}'' verdict of UC14: a document containing a provably false
theorem that contradicts its own sibling cannot be green. The honest
re-band is \AMBER{one false residual}: $R2$ is false; $R1$ (the
abutment identification) and $R3$ (the top-Walsh cofiber-LES argument)
were not re-audited and are left, respectively, at provisional and
``method-sound, soft spots noted''. The surviving honest statement of
UC13 Lemma~4.4.1 is its \emph{original} cautious \emph{surjection}
$\widetilde H^k(X(\Dx))_{\walsh S}\twoheadrightarrow\isotype
Sk(\Xn)$ --- consistent with $Y_1$, but not a proof of it (a surjection
onto a group does not show the group is zero).
\end{remark}

\subsubsection{The valid reduction}

The correct method --- the one UC14~$R3$ used for the sibling
top-Walsh statement --- is a cofiber long exact sequence with induction
on $n$. Carried out honestly for $Y_1$, it does not close $Y_1$, but it
\emph{validly reduces} it.

\begin{theorem}[Valid cofiber-LES reduction of $Y_1$]
\label{thm:y1-reduction}
Fix $n$, $S\subsetneq[n]$ with $|S|\le n-2$, $k\ge1$, and
$x\in[n]\setminus S$. Assume Conjecture~\ref{conj:uc10-1} for all
ground sets of size $<n$. The pair $(\Xn,\Xnm)$ has a
$(\Z/2)^n$-equivariant cohomology long exact sequence; projecting to
the $\walsh S$-isotype \textup(Maschke\textup) and substituting
$\iota_x^{\mathcal D\,*}=0$ \textup(Theorem~\ref{thm:trace-injectivity},
twisted form: the trace is null-homotopic\textup) collapses it to an
isomorphism
\[
  j^{*}\colon\;\isotype Sk(\cofib\iota_x)\;\xrightarrow{\ \cong\ }\;
  \isotype Sk(\Xn),\qquad \cofib\iota_x=\Xn/\Xnm.
\]
Hence, for $|S|\le n-2$, $Y_1$ holds \emph{if and only if}
$\isotype Sk(\cofib\iota_x)=0$ for all $k\ge1$.
\end{theorem}

\begin{proof}[Proof sketch]
The twisted bridge gives $\iota_x^{\mathcal D\,*}=0$ on
$\walsh S$-cohomology in every degree $k\ge1$ (the honest output of
Correction~\ref{corr:trace-exact}'s identity, \emph{fed into} the LES
rather than misread as vanishing). Substituting the zero map into the
isotype-projected LES makes the restriction map $j^{*}$ surjective with
kernel the image of the connecting map $\delta$, and $\delta$ injective;
the smaller-cube term $\isotype S{k-1}(\Xnm)$ vanishes by the inductive
hypothesis (it is a \emph{lower} isotype of a size-$(n-1)$ cube,
because $|S|\le n-2$ permits $S\subsetneq[n]\setminus\{x\}$ strictly).
The short exact sequence collapses to the stated isomorphism. The
$|S|=n-1$ sub-case is genuinely different: there $[n]\setminus S=\{x\}$
is forced, the smaller-cube term is the \emph{top} isotype of the
$(n-1)$-cube --- the nonzero sphere class --- and the argument requires
the separate $R3$-style machinery plus a sharp no-collision constraint.
\end{proof}

\begin{openresidual}[The $Y_1$ residual]
\label{res:y1}
$Y_1$ is reduced, by Theorem~\ref{thm:y1-reduction}, to the
\emph{vanishing of $\isotype Sk(\cofib\iota_x)$} --- the
$\walsh S$-isotype cohomology of the homotopy cofiber of the
matched-cylinder inclusion $X(\Dx)\hookrightarrow\Xn$. Stripping the
contractible cylinder, this is $\isotype Sk(\Xn/X(\Dx))$: the
cohomology of the Bousfield--Kan double complex of $\Cn$ indexed by the
bar strings that \emph{exit} the matched subcategory $\Dx$.
Equivalently, it is the reduced contribution of the non-contractible
\emph{comma categories} $\rho_x/(T,\G)$ of the trace functor $\rho_x$
--- the homotopy type of the ``stuck-family fibres'', the genuine
core of the program that has never been computed. \textbf{This is the
single open computation the cohomological closure of Frankl now rests
on.}
\end{openresidual}

\begin{remark}[Why the residual has no free vanishing]
\label{rem:no-free-vanishing}
Residual~\ref{res:y1} is not closeable by any of the program's
free mechanisms.
\emph{Not by family-by-family vanishing}: the cofiber
$\Xn/X(\Dx)$ is a hocolim, and at a bar string exiting $\Dx$ the
relative term is the \emph{entire} cochain group of a family, generically
nonzero --- termwise vanishing fails, exactly as in
Correction~\ref{corr:r2}.
\emph{Not by degree truncation}: $Y_1$'s range $1\le k\le n-2$ is the
middle range, with no dimension cap.
\emph{Not by the induction on $n$}: the inductive hypothesis controls
the smaller cube and the cylinder, not the cofiber, which has no
size-$(n-1)$ model.
\emph{Not by cofinality of $\rho_x$}: assuming $\rho_x$ homotopy-cofinal
is logically equivalent to $R2$, hence false
(Correction~\ref{corr:r2}). The residual must be \emph{computed}.
\end{remark}

\subsubsection{The Phase-1 sharpening: $Y_1$ as coefficient-blindness}

Theorem~\ref{thm:y1-reduction} and Open residual~\ref{res:y1} are the
standing of $Y_1$ after the cofiber-LES re-derivation. A subsequent
paper-and-pencil phase --- the $Y_1$-closure Phase~1 --- did \emph{not}
close the residual, but sharpened it. We record the sharpening, because it
reframes what a closure must do.

\begin{proposition}[The residual diagram is Walsh-index-independent]
\label{prop:y1-S-independent}
Fix $x$. The cofiber double complex $\cofib\iota_x=\Xn/X(\Dx)$ is
assembled from two pieces of data: the \emph{diagram shape} --- the
Bousfield--Kan bar strings of $\Cn$ that exit the matched subcategory
$\Dx$, with their bar differentials, together with the trace functor
$\rho_x$ and its comma categories --- and the \emph{coefficient system}
$G_S(\G):=C^*(X(\G);\Q)_{\walsh S}$ on $\mathcal C^{\cap}_{n-1,x}$. The
$\walsh S$-isotype cochain complex is
\[
  \BK^{p,q}\bigl(\Xn/X(\Dx)\bigr)_{\walsh S}\;=\;
  \bigoplus_{\text{bar strings exiting }\Dx}\;G_S(\rho_x\F_p),
\]
the pullback $\rho_x^{*}G_S$ restricted to the $\Dx$-exiting strings. The
diagram shape is \emph{independent of $S$}; only the coefficient system
$G_S$ depends on $S$.
\end{proposition}

\begin{proof}
The matched subcategory $\Dx$, the trace functor $\rho_x$, the property
``exits $\Dx$'', and the bar differentials are defined from $\Cn$ and the
choice of $x\in[n]\setminus S$ alone; none refers to $S$. The Walsh index
enters only through the isotype projection $(-)_{\walsh S}$ defining
$G_S$. The displayed identity is the cofiber quotient of the
Bousfield--Kan double complex (Definition~\ref{def:cubical-defect-complex})
intersected with the $\walsh S$-isotype: by UC14 Lemma~2.2.1 (the
cell-level Walsh-support lemma, sound at the single-family level) the
$\walsh S$-value $C^*(X(\F_p))_{\walsh S}$ depends only on the matched
part, hence only on the trace $\rho_x\F_p$.
\end{proof}

\begin{remark}[$Y_1$ as a coefficient-blindness statement]
\label{rem:y1-coefficient-blindness}
Proposition~\ref{prop:y1-S-independent} reframes $Y_1$. The residual must
\emph{vanish} for $|S|\le n-2$, yet at $S=[n]\setminus\{x\}$ it must carry
\emph{exactly} the sphere class $\sgn$ --- the genuine nonzero class of
Correction~\ref{corr:r2}. By Proposition~\ref{prop:y1-S-independent} this
contrast is \emph{not} a difference in the homotopy type of the
stuck-family fibres: that type is one fixed object, the same for every
$S$. It is a difference in which coefficient systems can \emph{see} it.
So $Y_1$, for $|S|\le n-2$, is precisely the assertion that \emph{every
lower-Walsh coefficient system $G_S$ with $|S|\le n-2$ is blind to the
fixed residual diagram} --- its reduced $G_S$-twisted cohomology vanishes
--- whereas the top-of-the-smaller-cube system $G_{[n]\setminus\{x\}}$ is
not blind (it sees $\sgn$). This is a representation-theoretic question
about the family $\{G_S\}$ of coefficient systems on one fixed diagram,
not a bare nerve computation: the comma-category nerves, being
$S$-independent, cannot by themselves account for the $|S|$-stratification
--- the coefficient system must be carried through. The residual is
sharpened, not closed.
\end{remark}

\begin{remark}[A foundational prerequisite]
\label{rem:y1-prerequisite}
Computing $\isotype Sk(\cofib\iota_x)$ --- by hand at the smallest open
instance $n=3$, $S=\{1\}$, $k=1$, or by machine, or in a faithful Lean
model --- requires the Bousfield--Kan model of $\Cn$ at \emph{cell} level:
the trace-induced structure map $X(\F|_T)\to X(\F)$, and the single-family
$\sigma_y$-action on $C^*(X(\F))$. Both are \emph{asserted} in the source
development (UC10~\S3.3; UC12~\S3.3) but not pinned to the precision a
verified computation needs. They are pinnable --- the trace map has a
canonical minimal-lift description, $A\mapsto\bigcap\{B\in\F:B\cap T=A\}$
--- but pinning them is a bounded prerequisite standing upstream of any
honest computation, independent validation, or formalisation of the
residual. A computation reported on un-pinned foundations would repeat the
failure mode of Corrections~\ref{corr:trace-exact}--\ref{corr:r2}.
\end{remark}

\begin{status}
$Y_1$ is \RED{open}. It is true-but-unproven: no counterexample is
known, the statement is plausibly true (it is the cube-side analogue of
Remark~\ref{rem:external-input}), but its only circulated proof is
invalid (Correction~\ref{corr:trace-exact}) and the lemma cited to
repair it is false (Correction~\ref{corr:r2}). The honest standing is
Theorem~\ref{thm:y1-reduction} plus Open residual~\ref{res:y1}: a
\emph{correct reduction} to one named homotopy-colimit computation,
estimated to be substantial (not a hypothesis-check). The $Y_1$-closure
Phase~1 did not close that computation: it re-verified the reduction,
sharpened the residual to the coefficient-blindness statement of
Remark~\ref{rem:y1-coefficient-blindness}, and isolated the foundational
prerequisite of Remark~\ref{rem:y1-prerequisite}; $Y_1$ stays
\RED{open}. Until the residual is computed, the two-part contradiction of
Section~\ref{ssec:closing} does not close, and the headline conjecture
is not proved.
\end{status}

\subsection{The machine-checked formalisation: honest state}
\label{ssec:lean-state}

A formalisation of the program exists in Lean~4 against
\texttt{mathlib}. We report its state honestly; it does not change any
mathematical verdict above.

\subsubsection{What is genuinely formalised}
\label{ssec:lean-genuine}

The following are real, machine-checked, and \texttt{sorry}-free:
the combinatorial bias arithmetic and the non-vanishing direction of
the obstruction (Theorem~\ref{thm:fact-one}); the
trace-injectivity / null-homotopy theorem
(Theorem~\ref{thm:trace-injectivity}); the $Y_2$ mechanism
(Theorem~\ref{thm:y2-mechanism}); and the concrete small cases of
Section~\ref{ssec:lean-unconditional}. The build is green.

\subsubsection{The headline theorem rests on one named axiom}
\label{ssec:axiom-is-frankl}

The formalised headline \texttt{Frankl\_Holds} --- the universal
statement, every nontrivial intersection-closed family has a rare
coordinate --- type-checks at every $n$, but its general branch routes
through a single named project axiom,
\texttt{case3\_ss\_obstruction\_paper\_axiom}. We must be precise about
what that axiom is.

\begin{proposition}[The project axiom is Frankl-in-disguise]
\label{prop:axiom-is-frankl}
The Lean development proves a theorem $X$ and assumes an axiom $Y'$,
both statements about the \emph{same} cohomology object
$\texttt{obstructionCohomClassChain}\,F$ under the hypothesis
$\texttt{IsCounterexample}\,F$:
\[
  X\ \textup{(proved)}:\ \texttt{IsCounterexample}\,F\Rightarrow
  \texttt{obstructionCohomClassChain}\,F\neq0,
\]
\[
  Y'\ \textup{(axiom)}:\ \texttt{IsCounterexample}\,F\Rightarrow
  \texttt{obstructionCohomClassChain}\,F=0.
\]
$X$ and $Y'$ are exact negations under the same hypothesis. In any
proof by contradiction, axiomatising one horn of a contradiction whose
other horn is proved is axiomatising the goal: $Y'$ together with the
proved $X$ yields $\forall F,\ \neg\,\texttt{IsCounterexample}\,F$
\textup(i.e.\ Frankl\textup) directly, and Frankl makes $Y'$ vacuously
true. Hence $Y'$ is propositionally equivalent to Frankl.
\end{proposition}

\begin{status}
\texttt{Frankl\_Holds} is \AMBER{axiom-dependent}, and the axiom is
\emph{not} a genuine sub-lemma --- it is Frankl itself, stated at the
chain level. This is not a framing accident: it is the mechanical
consequence of axiomatising one horn of a contradiction. A genuine
sub-lemma would have to be a statement about a \emph{different} object
than the one $X$ lives in --- the spectral-sequence abutment, or a
Walsh-isotype-refined chain group --- and that object does not exist in
the formalisation today (its construction is blocked by a
\texttt{mathlib} typeclass-instance diamond at the bicomplex level).
The formalisation discloses this honestly in its own axioms file. The
correct description of the artefact is therefore: \emph{the
combinatorial half is formalised; the homological half is assumed, and
as assumed it is --- given the formalised half --- equivalent to Frankl
itself.} It is not ``$X$ and $Y$ together prove Frankl''.
\end{status}

\subsubsection{The answer-steered isotype encodings}
\label{ssec:answer-steered}

Two further formalisation artefacts must not be cited as evidence.
First, the existing Lean ``formalisation of $Y_2$'' defines the
isotype projection as a hand-coded conditional
(\texttt{cechIsotypeProjection F x T := if T = \{x\} then \dots\ else 0})
--- it applies no group action and computes no projection, so proving
``the obstruction lands in level $1$'' from it is circular. Second, the
Lean ``formalisation of $Y_1$'' (\texttt{UC10\_lowerWalshVanishing})
proves a degree-$0$ identity of the form $(\text{operator})\,g=g$ whose
entire content is the Walsh involution $\chi^2=1$; it is a non-faithful
placeholder, with no $d$, no cocycle, no homology --- it neither
supports nor contradicts $Y_1$. Both are \emph{answer-steered}
encodings; the genuine, neutral $Y_2$ determination is
Theorem~\ref{thm:y2-mechanism}, and there is no faithful formalisation
of $Y_1$ (there is no valid proof to formalise).

\subsubsection{The genuinely unconditional slice}
\label{ssec:lean-unconditional}

The formalisation \emph{does} contain genuinely unconditional,
axiom-free machine-checked content: the concrete cases. For the full
powerset families on $[3]$ and on $[4]$ --- nontrivial
intersection-closed families, not degenerate --- the conclusion
``some coordinate is rare'' is verified directly, with the explicit
witness $x=0$ (where the bias is exactly $0$), computed by
decision procedure. These route through the unconditional concrete
slice \texttt{Frankl\_Holds\_concrete} (which takes a rare-coordinate
witness as a hypothesis) and \emph{bypass the obstruction bridge and
the project axiom entirely}; an audit of their axiom dependencies
confirms only the three \texttt{mathlib} foundational axioms.

\begin{status}
The honest unconditional machine-checked content is exactly this: the
full-powerset families at $n=3$ and $n=4$, and any family for which a
rare coordinate can be exhibited directly. The \emph{universal}
statement at $n=3$ and $n=4$ (every intersection-closed family on those
ground sets) is, in the present formalisation, routed through the
dispatch to the general branch and is therefore axiom-dependent. We
state this precisely to avoid overclaiming: ``$n=3$ and $n=4$ are
unconditional'' is true of the \emph{concrete full-powerset families},
not yet of the universally quantified statement, in the formalisation
as it stands. (Frankl is of course classically known for small ground
sets by other means~\cite{bruhn-schaudt}; that is not what the
formalisation encodes.)
\end{status}
